\documentclass[12pt,letterpaper]{article}

% ---- Packages ----
\usepackage[utf8]{inputenc}
\usepackage[T1]{fontenc}
\usepackage{amsmath,amssymb}
\usepackage{newtxtext,newtxmath}  % Times New Roman text + math (replaces mathptmx)
\usepackage[margin=1in]{geometry}
\usepackage{setspace}
\usepackage{titlesec}
\usepackage{enumitem}
\usepackage{xr-hyper}
\usepackage{hyperref}
\externaldocument{1 - General Overview}
\externaldocument{4 - Magnetic and Electric Fields}
\externaldocument{6 - Cosmology}
\externaldocument{7 - Quantum Mechanics}
\usepackage{microtype}
\usepackage{booktabs}

% ---- Section formatting ----
\titleformat{\section}{\large\bfseries\sffamily}{\thesection.}{0.5em}{}
\titleformat{\subsection}{\normalsize\bfseries\sffamily}{\thesubsection}{0.5em}{}
\titleformat{\subsubsection}{\normalsize\bfseries}{\thesubsubsection}{0.5em}{}

% ---- Spacing ----
\setstretch{1.15}
\setlength{\parskip}{0.5em}
\setlength{\parindent}{0em}

% ---- Custom commands ----
\newcommand{\rhoa}{\rho_{\!A}}
\newcommand{\Ka}{K_{\!A}}
\newcommand{\vin}{v_{\mathrm{in}}}
\newcommand{\rs}{r_s}
% \hbar is already defined in LaTeX; no redefinition needed

\begin{document}

\begin{center}
{\LARGE\bfseries Unifying Theory of Dynamic Aether Mechanics:\\
Gravity}\\[1em]
{\large Erik Otto}
\end{center}

\vspace{1em}

\begin{abstract}
Within the Dynamic Aether Mechanics framework, gravity arises as the residual
asymmetry of a transient binding cycle in which virtual particle--antiparticle pairs
form and annihilate continuously at every atom.  This paper develops the gravitational
mechanism, derives the gravitational constant $G$ from binding physics, and resolves
the vacuum catastrophe.  The framework is extended to gravitomagnetism (frame dragging,
confirmed by Gravity Probe~B), gravitational lensing (deriving the full
general-relativistic bending angle $4GM/(r_0 c^2)$ from two independent mechanisms),
and Cherenkov radiation.  Connections to mass--energy equivalence and the self-limiting
behavior of gravity are established.
\end{abstract}

\vspace{1em}

\bigskip
\noindent\textit{This paper is part of a series presenting the Unifying Theory of
Dynamic Aether Mechanics.  For the foundational concepts of Aether binding,
the three independent properties (density, pressure, temperature), and the
unification of force constants, see Paper~1: General Overview \cite{otto-overview}.}
\bigskip

\section{Gravity}
\label{sec:gravity}
% ============================================================

Gravity in this theory is a true wave---an oscillating pressure disturbance propagating
through Aether, generated by the transient bind--unbind cycling that occurs continuously
at every atom.\ The mechanism corresponds to the virtual particle--antiparticle fluctuations
described by quantum field theory, reinterpreted as the physical process of Aether
transiently binding and unbinding.


\subsection{The Mechanism}

The spinning vortex core of each electron acts as a rapidly moving boundary in the
sense of the Dynamic Casimir Effect \cite{wilson2011}: the circulation velocity near the
healing length ($\xi \approx 0.2$~fm) is sufficient to enhance vacuum fluctuations into
transient particle--antiparticle pairs---virtual pairs in the language of quantum field
theory.\ This is the same DCE mechanism that drives permanent binding
(Section~\ref{sec:emc2}), operating below the permanence threshold.\ Each binding event
consumes Aether from the local region, creating a brief low-pressure pulse.\ The pair is
unstable and almost immediately annihilates, restoring the consumed Aether and returning
the local pressure to normal.\ The vortex circulation then drives the next cycle.

Each cycle produces one pulse of the gravitational wave: a dip to low pressure followed
by a return to normal.\ At the extraordinarily high cycling frequencies implied by atomic-scale processes, this oscillation would be indistinguishable from a static attractive force
at any macroscopic measurement scale---consistent with our observation that gravity appears
to be a constant, static field.

The wave propagates outward at the speed of light, since light speed represents the
propagation speed of compression waves in Aether.\ This is consistent with the observed
speed of gravitational waves.


\subsection{Why Gravity Is Attractive}

The pressure wave produced by transient cycling is asymmetric: it oscillates between normal
pressure and low pressure, never exceeding normal pressure.\ The binding event consumes
Aether (low pressure), and the annihilation merely restores it (normal pressure).\ The time-averaged effect is therefore a net reduction in pressure near the source---a persistent
inward pull on surrounding matter.


\subsection{The Gravitational Constant and the Uncertainty Principle}

Each atom cycles independently, so the total gravitational wave output of a mass is
proportional to the number of atoms it contains---proportional to its mass $M$.\ The
gravitational constant $G$, in this framework, encodes properties of the transient cycling
process:

\begin{itemize}[leftmargin=2em]
  \item The mass of the transiently formed particle pair (how much Aether is consumed per cycle)
  \item The cycling frequency (how many bind--unbind events occur per second)
  \item The minimum vortex velocity at which transient binding is triggered
  \item The ratio of transient to permanent binding events (governed by the DCE permanence
    threshold)
\end{itemize}

The total force on a unit of mass in $m_2$ from all units in $m_1$ is $m_1 G$.\ Summed across
all units of $m_2$, the force at zero radius becomes $G m_1 m_2$.\ Since the pressure wave
dissipates radially, the force at distance $r$ becomes $Gm_1 m_2/r^2$---exactly as observed.

The cycling frequency and transient particle mass are not arbitrary.\ In quantum field
theory, virtual particle formation is governed by the Heisenberg uncertainty principle \cite{heisenberg1927}:
%
\begin{equation}
  \Delta E \cdot \Delta t \geq \frac{\hbar}{2}
\end{equation}
%
where $\Delta E$ is the energy of the virtual pair (related to its mass by $E = mc^2$) and
$\Delta t$ is the duration of its existence.\ More massive virtual pairs exist for shorter
durations.\ Less massive pairs can persist longer.\ This relationship constrains the cycling
parameters:

\begin{itemize}[leftmargin=2em]
  \item If the transient particle mass is large, it exists very briefly and cycles rapidly.
  \item If the transient particle mass is small, it exists longer and cycles more slowly.
\end{itemize}

In either case, the product of mass and cycling frequency is constrained by $\hbar$.\ The
Dynamic Casimir Effect further constrains the system by defining the vortex velocity at
which transient events become permanent.\ Together, these constraints mean $G$ may ultimately
be derivable from $\hbar$, $c$, and the specific properties of Aether (density, vortex
velocity threshold)---connecting gravity to quantum mechanics through fundamental constants.\ The
difficulty of reconciling gravity with quantum mechanics (the long-sought theory of quantum
gravity) may stem from not recognizing that they share a common mechanism: the transient
formation and annihilation of particles from a dynamic medium, governed by the same
uncertainty principle and DCE threshold conditions.

$G$ may not be truly constant.\ Since $G \propto \epsilon \, \rhoa \, \psi^2$ (as
derived in the Unification of Force Constants section), it could in principle vary through
two channels: the local Aether density $\rhoa$ and the asymmetry fraction $\epsilon$.\
However, the Bernoulli analysis shows that $\rhoa$ is essentially constant near
gravitational sources ($\Delta\rhoa/\rhoa = 0$ for free-fall inflow), and the cosmological
density variation is only $\sim$10$^{-4}$---negligible.\ The primary channel for $G$
variation is therefore through $\epsilon$, which depends on temperature: higher temperature
provides thermal energy that resists binding (analogous to how heating prevents
condensation), potentially reducing the fraction of each cycle's energy that escapes as
the gravitational wave.\ Cold environments (near matter, where thermal energy is locked
into mass formation) could have higher $\epsilon$ and therefore stronger gravitational
output, while hot environments (cosmic voids) could have lower $\epsilon$.


\subsection{Self-Limiting Behavior}

A potential concern is runaway feedback: if a transiently formed particle pair has mass,
and mass generates gravity, could transient particles trigger additional transient binding
in a cascading loop?

The Heisenberg uncertainty principle provides the primary self-limiting mechanism.\ A
virtual pair of energy $\Delta E$ can only exist for a time $\Delta t \leq \hbar / (2\Delta E)$.
The more massive the transient pair, the shorter its existence.\ A transient particle cannot
persist long enough to complete its own gravitational cycling---its lifetime is
constrained by the uncertainty principle to be far shorter than the period of any
gravitational wave it could produce.\ The cascade is cut off by the same quantum mechanical
constraint that governs the cycling in the first place.

The Dynamic Casimir Effect reinforces this: for a transient particle to trigger further
permanent binding, it would need to act as a ``moving boundary'' with sufficient velocity
to exceed the DCE threshold.\ But a transient particle exists for too brief a duration and
has no persistent spin vortex---it lacks the sustained, directional motion required to
function as a DCE boundary.\ Only permanent particles with stable electron spin can sustain
the vortex dynamics needed for continuous cycling.

Additional self-limiting factors include:

\begin{itemize}[leftmargin=2em]
  \item The annihilation of the transient pair produces a restoring pressure pulse that
    disrupts the circulation pattern needed to sustain further binding nearby.

  \item The DCE threshold acts as a natural regulator: only Aether circulated by a
    persistent vortex (a permanent electron) undergoes cycling.
    A fleeting fluctuation from a transient pair does not sustain the boundary
    motion long enough to trigger secondary binding.
\end{itemize}


\subsection{Connection to \texorpdfstring{$E = mc^2$}{E = mc²}}
\label{sec:emc2}

The transient binding mechanism connects directly to mass--energy equivalence.\ Each bind--unbind cycle momentarily converts energy into bound mass and back.\ The energy recycled per
cycle is $mc^2$ of the transiently formed particle pair (as derived in the Special Relativity
section below).\ In QFT terms, the energy is ``borrowed'' from the vacuum for duration
$\Delta t$ and returned upon annihilation.

The gravitational wave's energy is a tiny fraction of $mc^2$---only the net outward-propagating pressure disturbance, not the full binding energy.\ This is why gravity is so
extraordinarily weak compared to other forces: the vast majority of each cycle's energy
is recycled locally, and only a minuscule asymmetry (the difference between the low-pressure binding pulse and the normal-pressure restoration) propagates as the gravitational
wave.

When the DCE threshold is exceeded and a binding event becomes permanent, the full $mc^2$
of energy is locked into the resulting particle rather than being returned.\ This is the
energy cost of mass formation, paid by the local environment (the vortex kinetic energy
and surrounding Aether thermal energy).


\subsection{The Vacuum Catastrophe}

The identification of transient binding with virtual particle fluctuations offers a
perspective on one of the largest unsolved problems in physics: the vacuum catastrophe.

When physicists use quantum field theory to calculate the energy density of vacuum
fluctuations, they obtain a value roughly $10^{120}$ times larger than the observed
cosmological constant (the measured energy density of empty space).\ This is often called
the worst prediction in the history of physics.

This theory suggests a resolution.\ The enormous energy of vacuum fluctuations does exist---it is the total energy being cycled through transient binding events at every point in
space.\ But nearly all of this energy is recycled locally.\ Each binding event invests $mc^2$
and recovers $mc^2$ upon annihilation.\ The net energy that propagates outward as the
gravitational wave is only the tiny asymmetry of the cycle---the difference between the
low-pressure dip and the normal-pressure restoration.

The vacuum catastrophe arises from counting the total energy of fluctuations.\ The
physically observable effect (gravity, and by extension the cosmological constant) comes
only from the residual asymmetry.\ The ratio between the two---the fraction of vacuum
energy that actually propagates rather than being recycled---could account for the
enormous discrepancy.\ The vacuum energy is real, but its gravitational effect is not its
total magnitude; it is only the unrecycled remainder.


% ============================================================


\subsection{Gravitomagnetism: Frame Dragging from Circulatory Aether Flow}
\label{sec:gravitomagnetism}

The $v^2/c^2$ universality established in Section~\ref{sec:force-constants} has a direct
gravitational consequence.  In electromagnetism, a moving charge sweeps its flow vector
(the electric field) into angular structure (the magnetic field), with
$F_{\mathrm{mag}} = F_{\mathrm{elec}} \times v_1 v_2 / c^2$.\footnote{Paper~4~\cite{otto-em}
identifies the electromagnetic ``flow vector'' concretely as the coherent
axial-oscillation wave field of the bound core ring; the gravitational analog here uses
the same flow-vector/flow-angle dichotomy at the macroscopic level.}
The same mechanism applies to any source of Aether perturbation that is set into
motion---including gravity.

A mass element moving at velocity $\mathbf{v}$ drags its locally bound Aether,
creating angular flow structure that stands in the same relation to the
gravitoelectric flow vector $\mathbf{g}$ as the magnetic flow angle $\mathbf{B}$
stands to the electric flow vector $\mathbf{E}$:
\begin{equation}
  \mathbf{B}_g \;=\; -\,\frac{1}{c^2}\,\mathbf{v}\times\mathbf{g}\,.
  \label{eq:Bg-from-v}
\end{equation}
This is the gravitomagnetic field.  No new parameters are introduced; the coupling
is fixed entirely by the ratio $v^2 \rhoa / K_A = v^2/c^2$.

For a body with angular momentum $\mathbf{J}$, integration over the
mass--current distribution yields a dipole gravitomagnetic field identical in
structure to the magnetic dipole field of a current loop:
\begin{equation}
  \mathbf{B}_g(\mathbf{r})
  \;=\; \frac{G}{c^2\,r^3}\,
        \bigl[3(\mathbf{J}\cdot\hat{\mathbf{r}})\,\hat{\mathbf{r}}
              - \mathbf{J}\bigr]\,.
  \label{eq:Bg-dipole}
\end{equation}
A gyroscope with spin angular momentum $\mathbf{S}$ placed in this field
precesses at the Larmor rate
\begin{equation}
  \boldsymbol{\Omega}_{\mathrm{FD}}
  \;=\; \frac{1}{2}\,\mathbf{B}_g\,,
  \label{eq:FD-precession}
\end{equation}
where the factor of $\tfrac{1}{2}$ is the standard gyroscopic coupling
(spin precesses at half the imposed rotation rate).

\paragraph{Quantitative prediction for Gravity Probe~B.}
The Gravity Probe~B (GP-B) satellite placed four cryogenic gyroscopes in a
polar orbit at altitude $h = 642$~km (orbital radius $R = 7.013\times10^6$~m)
above the spinning Earth ($J_\oplus = I_\oplus\,\omega_\oplus
= 5.86\times10^{33}$~kg\,m$^2$\,s$^{-1}$).
The gyroscope spin axes were aligned with the guide star IM~Pegasi
(declination $\delta = 16.84^\circ$).
Averaging over the polar orbit, the frame-dragging precession in the
west--east direction is
\begin{equation}
  \Omega_{\mathrm{FD}}
  \;=\; \frac{G\,J_\oplus}{2\,c^2\,R^3}\,\cos\delta
  \;=\; 39.3~\text{mas\,yr}^{-1}\,,
  \label{eq:FD-GPB}
\end{equation}
where $1~\text{mas} = 4.85\times10^{-9}$~rad.
This is a \emph{no-fit derivation}: it follows entirely from the $v^2/c^2$
universality with zero adjustable parameters.

The GP-B experiment measured $\Omega_{\mathrm{FD}} = -37.2 \pm 7.2$~mas\,yr$^{-1}$,
in agreement with the prediction within $0.3\sigma$.
General relativity predicts $-39.2$~mas\,yr$^{-1}$ from the weak-field limit of
the Kerr metric---the same formula, confirming that the Aether mechanism
reproduces the Lense--Thirring effect exactly.
The LAGEOS satellite tracking experiments independently confirmed frame dragging
to approximately $10\%$ precision via orbital node precession.

\paragraph{Physical interpretation.}
The Aether theory provides a concrete mechanical picture for frame dragging.
A rotating mass continuously drags its permanently bound Aether into
co-rotation.  This angular flow structure is structurally identical to the
pattern that produces magnetism from moving charges---in both cases, motion
converts a flow vector into a flow angle---differing only in the coupling
mechanism (mass binding versus charge flow) and therefore in overall strength
(suppressed by $\varepsilon \sim 10^{-43}$ relative to EM).

This unifies the concept of ``frame dragging'' discussed in Section~\ref{sec:black-holes}
(Black Holes) with the magnetic force mechanism of Section~\ref{sec:magnetism} (Magnetism):
both are circulatory Aether flow from a moving source, governed by the
universal $v^2/c^2$ ratio.  For black holes, the same mechanism drives
relativistic jet collimation, where the central rapidly rotating mass
creates gravitomagnetic fields strong enough to redirect infalling matter
along the polar axis.

\paragraph{Fourth $v^2/c^2$ phenomenon.}
Gravitomagnetism extends the $v^2/c^2$ universality from three confirmed
phenomena to four:
\begin{enumerate}[label=(\alph*)]
  \item \textbf{Time dilation}: $\gamma = 1/\sqrt{1 - v^2/c^2}$
        --- geometric path-length effect,
  \item \textbf{Magnetic force}: $F_{\mathrm{mag}} = F_{\mathrm{elec}}
        \times v_1 v_2/c^2$ --- flow angle from charge motion,
  \item \textbf{Relativistic energy}: $E = \gamma mc^2$
        --- compression energy,
  \item \textbf{Gravitomagnetism}: $B_g = g\times v/c^2$
        --- flow angle from mass motion; confirmed by GP-B.
\end{enumerate}
All four are manifestations of the single dimensionless ratio
$v^2 \rhoa / \Ka = v^2/c^2$: kinetic energy density divided by the
elastic capacity of the Aether medium.


\section{Local Effects: Cherenkov Radiation and Gravitational Lensing}
\label{sec:local-effects}
% ============================================================

These phenomena are explained by the Aether inflow environment near matter.

\textbf{Cherenkov radiation} is a flash of light that occurs when a charged particle
traveling at near light speed enters a material medium with a lower speed of light, such
as water or glass.\ In such materials, the bound atomic structure modifies the Aether's
wave propagation properties: the effective wave speed is reduced to $c/n$, where $n$
is the material's refractive index.\ When a fast particle enters this region and exceeds
the local wave propagation speed, it produces a compression shockwave---a ``lucid
boom''---that either is the observed light itself or causes the release of photons by
disturbing nearby electrons.\ Note that this is a property of \emph{material} media (bound
matter), not of the gravitational inflow: the free Aether maintains $c = \sqrt{\Ka/\rhoa}$
throughout the gravitational field, as shown by the Bernoulli analysis below.


\subsection{Gravitational Lensing}
\label{sec:lensing}

Gravitational lensing is explained by the acoustic metric of the Aether inflow.

\subsubsection{The Acoustic Metric}

Light propagating through a flowing medium follows geodesics of the \emph{acoustic
metric}---the effective spacetime geometry experienced by waves in a moving medium.\
For the gravitational inflow at $\vin(r) = \sqrt{2GM/r}$ with uniform wave speed~$c$,
the acoustic metric takes the Painlev\'e--Gullstrand form:
%
\begin{equation}
  ds^2 = \bigl(c^2 - \vin^2\bigr)\,dt^2
    - 2\,\vin\,dr\,dt - dr^2 - r^2\,d\Omega^2
\end{equation}
%
This is mathematically equivalent to the Schwarzschild metric of general relativity,
expressed in ``river'' coordinates where the Aether flows inward at escape velocity.\
All predictions of the Schwarzschild geometry---including gravitational lensing, time
dilation, and the event horizon---follow directly from this acoustic metric without
invoking spacetime curvature.

\subsubsection{Deflection Angle}

The total bending angle for a light ray passing at closest approach $r_0$ from a mass
$M$ follows from the geodesics of the acoustic metric:
%
\begin{equation}
  \theta_{\mathrm{total}} = \frac{4GM}{r_0 \, c^2}
  \label{eq:lensing}
\end{equation}
%
This can be decomposed into two contributions of equal magnitude:
$2GM/(r_0 c^2)$ from transverse advection by the inflow (the light beam is
carried sideways by the flowing medium) and $2GM/(r_0 c^2)$ from the
time-delay asymmetry (the beam spends more coordinate time in the near-side
flow than the far-side, creating a net deflection).\ Both contributions arise
from the same inflow---they are aspects of a single physical mechanism, not
two independent effects.

This matches the observed gravitational deflection confirmed since Eddington's 1919
solar eclipse measurements and subsequently verified to better than 0.01\% precision by
radio interferometry \cite{bertotti2003}.\ The Aether model reproduces the full
general-relativistic bending angle from the acoustic metric of a flowing medium, without
invoking spacetime curvature.

\subsubsection{The Event Horizon}

The inflow velocity $\vin(r) = \sqrt{2GM/r}$ increases as $r$ decreases.\ At the Schwarzschild
radius $\rs = 2GM/c^2$, the inflow reaches the speed of light:
%
\begin{equation}
  \vin(\rs) = \sqrt{\frac{2GM}{2GM/c^2}} = c
\end{equation}
%
At this radius, no outward-propagating wave can escape---not because spacetime geometry
prevents it, but because the medium itself is flowing inward at the wave propagation speed.
This is directly analogous to a sonic horizon in fluid dynamics, where a supersonic flow
prevents upstream sound propagation.

Inside the Schwarzschild radius, the inflow exceeds $c$.\ The transient binding cycle still
operates locally (pair formation and annihilation do not require outward propagation), but
the restoring pressure pulses from pair annihilation can no longer propagate outward.\ The
energy from annihilation is carried further inward by the flow, potentially leading to net
Aether accumulation at the center.

\subsubsection{Steady-State Conservation}
\label{sec:steady-state}

The Aether inflow at escape velocity raises an immediate question: if Aether continuously
flows inward toward mass, where does it go?\ Does the mass grow from accumulating Aether?

The answer lies in the transient binding cycle that the theory already describes.\ The
inflow is not a separate process from gravity---it \emph{is} gravity, viewed from the fluid-dynamics perspective:

\begin{enumerate}[leftmargin=2em]
  \item Electron vortices draw Aether inward (creating the inflow).
  \item The vortex core acts as a rapidly moving boundary whose circulation velocity
    exceeds the DCE threshold, producing transient particle--antiparticle pairs
    (consuming Aether from the local region).
  \item The pair annihilates almost immediately, restoring the consumed Aether (returning it
    to the medium).
  \item The restored Aether is available to flow inward again, completing the cycle.
\end{enumerate}

The net Aether consumption is essentially zero---only the tiny asymmetry $\epsilon$
propagates outward as the gravitational wave.\ The inflow velocity $\vin(r) = \sqrt{2GM/r}$
represents the bulk drift velocity of a steady-state circulation pattern, not a one-way
accumulation.\ The kinetic energy of the inflow at radius $r$ is $\tfrac{1}{2}\rhoa \vin^2 =
\rhoa GM/r$, exactly equal to the gravitational potential energy density---consistent
with free-fall into a gravitational potential well with full energy conservation.

This picture connects to the vacuum catastrophe: the enormous energy of vacuum
fluctuations (approximately $10^{120}$ times the cosmological constant) is the total cycling
energy being invested and recovered at every point in space.\ The gravitational effect
comes only from the tiny unrecycled asymmetry.\ The Aether inflow is the macroscopic
manifestation of this microscopic recycling process.

This theory also predicts that gravitational waves propagate at the speed of light, since
light speed represents the propagation speed of compression waves in Aether.


% ============================================================

% ============================================================
\section{Summary of Theoretical Properties}
% ============================================================

\subsection{Properties of Gravity}

\begin{itemize}[leftmargin=2em]
  \item An oscillating pressure wave (normal $\rightarrow$ low $\rightarrow$ normal) produced by the transient
    bind--unbind cycling of Aether at every atom, corresponding to virtual particle--antiparticle fluctuations in quantum field theory.
  \item Triggered by the electron vortex core acting as a rapidly moving boundary---the
    DCE mechanism operating below the permanence threshold.
  \item At macroscopic scales, the extraordinarily high cycling frequency makes this
    oscillation indistinguishable from a static attractive field.
  \item Propagates at the speed of light.
  \item Requires no net consumption of Aether and no ejection of particles.\ Energy is
    recycled locally; only the pressure wave propagates outward.
  \item The weakness of gravity relative to other forces reflects the fact that nearly all
    energy in each cycle is recycled, with only a tiny asymmetry (the difference between
    the low-pressure dip and the normal-pressure restoration) propagating as the wave.
  \item The gravitational constant $G$ encodes the transient particle mass, cycling frequency,
    vortex velocity threshold, and transient/permanent ratio (DCE threshold), all constrained by
    the Heisenberg uncertainty principle ($\Delta E \cdot \Delta t \geq \hbar/2$) and the DCE
    mathematics.\ $G$ is predicted to be derivable from the Coulomb constant $k$, the speed of
    light $c$, and the asymmetry fraction $\epsilon$ of the bind--unbind cycle---reducing all
    force constants to properties of the same Aether medium.
  \item The Aether inflow velocity at radius $r$ is $\vin(r) = \sqrt{2GM/r}$, equal to the
    Newtonian escape velocity.\ At the Schwarzschild radius $\rs = 2GM/c^2$, the inflow
    reaches $c$, producing a sonic horizon (event horizon) beyond which no
    outward-propagating wave can escape.\ The inflow is a steady-state
    circulation: Aether flows in, participates in transient binding, and is
    returned upon pair annihilation.\ In the two-fluid model
    (Paper~6~\cite{otto-cosmology}), the gravitational inflow is the
    center-of-mass velocity of both the superfluid and normal components
    flowing inward together; the bind--unbind cycle operates on the
    superfluid condensate, which constitutes $\sim$100\% of the Aether near
    matter where $T \ll T_c$.
  \item Self-limiting feedback, governed by the uncertainty principle and the DCE requirement
    for sustained vortex motion, prevents runaway cascading of transient binding events.
  \item \textbf{Gravitomagnetism.}
    A rotating mass drags bound Aether into co-rotation, creating angular
    flow structure $\mathbf{B}_g = -(1/c^2)\,\mathbf{v}\times\mathbf{g}$---the
    gravitational flow angle, analogous to the electromagnetic flow angle, governed
    by the same $v^2/c^2$ ratio.  Confirmed by GP-B (frame-dragging precession
    $39.3$ vs.\ measured $37.2 \pm 7.2$~mas\,yr$^{-1}$).
\end{itemize}



% ============================================================
\section{Experimental Support and Connections to Established Physics}
% ============================================================

The following experimentally confirmed phenomena support or align with the aspects of the theory presented in this paper:

\begin{enumerate}[leftmargin=2em]
  \item \textbf{Cherenkov radiation.} Particles exceeding the local speed of light in a material
    medium produce light---consistent with this theory's prediction that bound matter
    modifies the Aether's wave propagation properties (reducing the effective wave speed
    to $c/n$), creating conditions for a compression shockwave.

  \item \textbf{Gravitational wave speed.} LIGO's detection of gravitational waves \cite{ligo2016}, and
    the near-simultaneous arrival of gravitational and electromagnetic signals from the
    neutron-star merger GW170817 \cite{gw170817}, confirmed they
    propagate at the speed of light to within $|{\Delta c}/{c}| < 10^{-15}$, consistent
    with this theory's prediction that gravitational waves are compression waves in the
    same medium that carries light.

  \item \textbf{Gravitational lensing bending angle.} The observed deflection of light by massive
    objects is $4GM/(r_0 c^2)$, confirmed since Eddington's 1919 solar eclipse \cite{eddington1920} and verified
    to better than 0.01\% by radio tracking \cite{bertotti2003}.\ This theory derives this exact result from
    the acoustic metric of the Aether inflow at escape velocity $\vin = \sqrt{2GM/r}$,
    which takes the Painlev\'e--Gullstrand form and is exactly equivalent to the
    Schwarzschild geometry of general relativity.

  \item \textbf{Gravitomagnetic frame dragging (GP-B, 2011; LAGEOS, 2004).}
    The $v^2/c^2$ universality applied to gravity predicts a frame-dragging
    precession of $39.3$~mas\,yr$^{-1}$ for the GP-B orbit
    (Eq.~\ref{eq:FD-GPB}), with no adjustable parameters.
    GP-B measured $-37.2 \pm 7.2$~mas\,yr$^{-1}$, consistent within
    $0.3\sigma$.  LAGEOS confirmed frame dragging independently
    to $\sim\!10\%$ precision~\cite{everitt2011,Ciufolini2004}.

  \item \textbf{Null results for static electromagnetic--gravitational coupling.}\
    Hathaway and Williams \cite{hathaway2021} found no weight change for
    conducting objects at potentials up to 200~kV (null to within 2~g), and
    Tajmar, K\"ossling, and Neunzig \cite{tajmar2024} tested capacitors,
    solenoids, and tunneling-current devices on high-resolution balances in
    vacuum, finding no anomalous forces down to $\sim$3~$\mu$N.\ The Aether
    model predicts exactly this null result: static electric fields are the
    time-averaged coherent axial-oscillation wave field of the bound core ring
    (Paper~4~\cite{otto-em}), a net-zero-flow oscillation that does not
    involve the transient bind--unbind cycling that generates gravity.\
    Only time-varying disruptions of the cycling process can modify
    gravitational output.


\end{enumerate}



% ============================================================
\section{Open Questions and Testable Predictions}
% ============================================================

The following areas relevant to this paper remain underdeveloped and represent opportunities for further refinement or falsification:

\begin{enumerate}[leftmargin=2em]
  \item \textbf{Deriving $G$ from fundamental constants.} If transient binding corresponds to virtual
    particle fluctuations governed by the Heisenberg uncertainty principle, and the
    transient/permanent ratio is governed by the DCE threshold, it may be possible to
    derive the gravitational constant from $\hbar$, $c$, the DCE threshold velocity, and
    the properties of Aether.\ A successful derivation would be strong evidence for this
    theory and would represent a unification of gravity and quantum mechanics.

  \item \textbf{The vacuum catastrophe ratio.} This theory predicts that the ratio of total vacuum
    fluctuation energy to the observed cosmological constant should equal the ratio of
    total cycling energy to the asymmetric residual that propagates as gravity.\ If this
    ratio can be calculated from the theory's parameters, it could resolve the vacuum
    catastrophe---one of the largest unsolved problems in physics.

  \item \textbf{Gravitational wave frequency.} If gravity is an oscillating wave rather than a static
    field, it has a characteristic frequency determined by the cycling rate.\ The
    uncertainty principle constrains this frequency.\ If it is finite (however high),
    it is in principle detectable.\ Identifying this frequency would be a strong
    confirmation of the theory.

  \item \textbf{Variation of $G$.} If $G$ depends on local Aether conditions (density, temperature),
    precision measurements of $G$ in different environments could reveal deviations from
    constancy.\ Such deviations, if found, would support this theory over standard models
    where $G$ is a universal constant.

  \item \textbf{Steady-state Aether inflow and interior solutions.} The gravitational inflow velocity
    $\vin(r) = \sqrt{2GM/r}$ reaches $c$ at the Schwarzschild radius, producing a sonic horizon
    from which outward-propagating waves cannot escape.\ Inside this radius, the transient
    binding cycle still operates but the restoring pressure pulses are carried further
    inward rather than propagating outward, potentially leading to net Aether accumulation
    at the center.\ The conservation analysis (steady-state recycling at the macroscopic
    level) should be extended to include the gravitational wave energy loss rate $\epsilon$
    and its long-term effect on mass evolution.

  \item \textbf{Self-consistent strong-field solution.} The weak-field inflow velocity
    $\vin(r) = \sqrt{2GM/r}$ is a first-order approximation valid when
    $2GM/(rc^2) \ll 1$.\ Near a black hole, where $\vin$ approaches~$c$,
    the full solution requires solving the inflow velocity profile, the
    acoustic metric, and the bind--unbind cycling rate self-consistently
    through the continuity and momentum equations (with the Aether's
    logarithmic equation of state).\ The Bernoulli analysis guarantees
    $\Delta\rhoa/\rhoa = 0$ in the weak field; the strong-field regime
    must determine whether this holds or whether relativistic corrections
    introduce density perturbations.\ This self-consistent solution would
    determine: (a)~the exact location of the sonic horizon, (b)~the precise
    photon sphere radius, (c)~whether the Aether acoustic metric reproduces the
    Schwarzschild metric in the strong-field regime, and (d)~the exact Hawking
    temperature coefficient.

  \item \textbf{Gravitational wave merger waveforms.} LIGO and Virgo \cite{ligo2016} have detected dozens of binary
    black hole mergers whose gravitational waveforms match general-relativistic numerical
    simulations with extraordinary precision, including the inspiral chirp, the merger
    peak, and the ringdown quasi-normal modes.\ The Aether model predicts gravitational
    waves propagate at $c$ (confirmed) and that they arise from compression waves in the
    medium.\ Reproducing the detailed merger waveforms---particularly the nonlinear merger
    and ringdown phases---from Aether dynamics would be a stringent test.\ This likely
    requires numerical simulation of colliding Aether inflow patterns.


  \item \textbf{Electromagnetic control of the asymmetry fraction
    $\epsilon$.}\ The gravitational constant $G \propto \epsilon\,\rhoa\,\psi^2$
    depends on the asymmetry fraction of the bind--unbind cycle.\ Can
    $\epsilon$ be modified by external electromagnetic fields?\ The null
    results for static fields \cite{hathaway2021,tajmar2024} constrain the
    coupling, but time-varying fields that resonantly drive the binding
    cycle have not been tested.\ Determining whether $\epsilon$ is a
    fixed property of the medium or a dynamically controllable parameter
    would have profound implications for the relationship between gravity
    and electromagnetism.

  \item \textbf{Cooper pair suppression of the gravitational cycling rate.}\
    In a superconductor, electrons form Cooper pairs whose opposite-spin
    vortices cancel all net Aether flow.\ Does this cancellation extend to
    the transient bind--unbind cycle?\ If Cooper pairing suppresses the
    cycling rate or modifies $\epsilon$ within the condensate, a
    superconductor should exhibit a measurable difference in gravitational
    coupling compared to the same material above its critical temperature.\
    Precision measurements of the gravitational attraction between a test
    mass and a sample cycled above and below $T_c$ could test this
    prediction.



\end{enumerate}



% ============================================================
\section{Proposed Experiments}
% ============================================================

The following experiments target specific predictions relevant to this paper:

\begin{enumerate}[leftmargin=2em]
  \item \textbf{Weight measurement above a driven superconductor with AC fields.}\
    The theory predicts that static electromagnetic fields cannot modify gravity
    (confirmed by null results \cite{hathaway2021,tajmar2024}), but that
    \emph{time-varying} fields applied to a superconductor could coherently
    drive the transient bind--unbind cycle, amplifying the asymmetry fraction
    $\epsilon$ and producing a measurable weight change in a test mass above the
    device.\ The proposed setup consists of a large-diameter YBCO or niobium
    superconductor subjected to alternating magnetic fields at frequencies from
    kHz to GHz, with a precision balance measuring test mass weight changes under
    vacuum.\ The theory predicts: (a)~no effect from static fields (already confirmed),
    (b)~a frequency-dependent onset when the driving field couples to the
    bind--unbind cycling rate, and (c)~the effect should vanish above the
    superconducting critical temperature, since Cooper pair coherence is
    required to make the cycling collective rather than random.

  \item \textbf{Refractive-index gradient as a gravitational potential.}\
    In the Aether model, a material with refractive index $n$ increases the
    effective Aether density by $n^2$ within its volume, reducing the local
    wave speed to $c/n$.\ (This is distinct from gravitational fields, where
    Bernoulli gives $\Delta\rhoa = 0$.)\ A metamaterial with
    an engineered refractive-index gradient therefore creates a gradient in
    local $c$, which is physically identical to a gravitational potential
    gradient.\ The proposed experiment uses a metamaterial slab with a
    monotonic refractive-index gradient from $n = 1$ to $n \approx 2$ over
    a distance of $\sim$10~cm, and measures clock rates (via precision
    atomic spectroscopy or optical frequency comparison) at different
    positions within the gradient.\ The theory predicts a fractional frequency
    shift of $\Delta f / f = (n_2 - n_1)/n_2$ between the two
    positions---since clock rates scale with the local wave speed
    $c/n$, the fractional shift is simply the ratio of local
    propagation speeds.

  \item \textbf{Asymmetric waveform gravity generation.}\
    The Aether model predicts gravity arises from an asymmetric pressure wave:
    the bind event produces a low-pressure pulse, and the unbind event partially
    restores pressure, but with a net deficit (the asymmetry fraction $\epsilon$).
    An artificial gravity source could in principle be constructed by driving
    an asymmetric electromagnetic waveform---fast compression followed by
    slow rarefaction---through a medium capable of coupling the electromagnetic
    oscillation to local Aether density changes.\ A superconductor is the natural
    candidate, since its Cooper pair condensate couples coherently to the medium.\
    The proposed experiment subjects a superconductor to a sawtooth-wave magnetic
    field (sharp rise, slow fall) and measures whether the asymmetric driving
    produces a net force on nearby test masses, compared to a symmetric sinusoidal
    drive at the same frequency and amplitude (which the theory predicts should
    produce no net effect).

  \item \textbf{Casimir cavity gravitational test.}\
    The dynamical Casimir effect has been confirmed in superconducting circuits
    \cite{wilson2011}, demonstrating that accelerating boundaries convert vacuum
    fluctuations into real particles.\ This is the microscopic mechanism by which
    the theory proposes mass forms from Aether.\ The proposed experiment measures
    whether a Casimir cavity undergoing rapid, asymmetric mechanical oscillation
    produces a detectable gravitational signal---either a static weight change
    or an oscillating force on a nearby test mass.\ The theory predicts that
    symmetric oscillation produces no net gravitational output (equal binding
    and unbinding), while asymmetric oscillation (where the boundary spends more
    time accelerating in one direction) should produce a net gravitational impulse
    proportional to the asymmetry and the DCE photon production rate.

  \item \textbf{Superfluid helium analog of gravitational time dilation.}\
    The 2024 observation of Kerr black hole geometry in a superfluid ${}^4$He
    vortex \cite{svancara2024} demonstrated that quantum superfluids reproduce
    gravitational spacetime structure.\ The next step is to measure the analog of
    \emph{gravitational time dilation}: whether wave propagation through the
    radial inflow of a draining vortex exhibits time dilation consistent with the
    Painlev\'{e}--Gullstrand acoustic metric, where a static observer in the
    inflow experiences $\delta t/t = \vin^2/(2c_s^2)$.\ Existing superfluid
    vortex facilities could perform this measurement by comparing wave arrival
    times at different radial distances from a stabilized giant vortex.\
    Agreement would confirm that gravitational time dilation arises from the
    inflow velocity field of a superfluid medium.

  \item \textbf{Frequency dependence of the gravitational constant.}\
    If gravity arises from an oscillating bind--unbind cycle rather than a static
    field, measurements of $G$ at different timescales could reveal the cycling
    frequency.\ Torsion balance measurements of $G$ (which operate at periods
    of minutes) should agree with atomic interferometry measurements (which
    operate at millisecond timescales).\ However, at frequencies approaching the
    bind--unbind cycling rate, the effective $G$ should begin to decrease as the
    measurement period becomes shorter than the cycling period.\ Comparing
    precision measurements of $G$ across the widest possible range of measurement
    timescales---from torsion balances to atom interferometers to high-frequency
    resonators---could reveal the characteristic frequency of the gravitational
    cycle.

\end{enumerate}

\noindent\textit{For all proposed experiments, see Paper~1: General Overview \cite{otto-overview}.}


% ============================================================
\section{Mathematical Summary}
% ============================================================

This section collects the key equations relevant to the topics covered in this paper.
\noindent\textit{For the complete mathematical summary, see Paper~1: General Overview \cite{otto-overview}.}

\subsection{Fundamental Aether Parameters}
% --------------------------------------------------

The theory posits four microscopic parameters from which all macroscopic
phenomena derive:
%
\begin{center}
\begin{tabular}{@{}llll@{}}
\hline
\textbf{Symbol} & \textbf{Name} & \textbf{Constrained range} & \textbf{Primary constraint} \\
\hline
$\rhoa$  & Ambient density     & $10^{11}$--$10^{12}$~kg/m$^3$         & Electron vortex energy \\
$\Ka$    & Bulk modulus         & $10^{28}$--$10^{29}$~J/m$^3$          & $\Ka = \rhoa\,c^2$ \\
$m_A$    & Aether quantum mass  & $500$--$850$~MeV/$c^2$                & Flux tube width / glueball mass \\
$\xi$    & Healing length       & $0.17$--$0.28$~fm                     & Bridge diameter (lattice QCD) \\
\hline
\end{tabular}
\end{center}

These four are not independent.\ Two defining relations connect them:
%
\begin{align}
  c &= \sqrt{\Ka / \rhoa}
    \label{eq:sum-speed-of-light} \\
  \xi &= \frac{\hbar}{m_A\,c\,\sqrt{2}}
    \label{eq:sum-healing-length}
\end{align}
%
leaving two free Aether parameters (e.g., $\rhoa$ and $m_A$) from which all
others follow.\ (A third parameter, $\epsilon$, enters through gravity; see below.)


% --------------------------------------------------
\subsection{Equation of State}
% --------------------------------------------------

The Aether obeys a logarithmic equation of state, the unique form that produces
a density-independent bulk modulus (Section~\ref{sec:qss}):
%
\begin{equation}
  P = \Ka \ln\!\Bigl(\frac{\rho}{\rho_0}\Bigr) + P_0
  \label{eq:sum-eos}
\end{equation}
%
Key consequence: the speed of sound $c_s = \sqrt{\Ka/\rho}$ varies with
density.\ At ambient density, $c_s = c$.\ Note: in gravitational fields,
Bernoulli's equation for free-fall inflow gives $\Delta\rho/\rhoa = 0$
exactly, so $c$ does not vary near gravitating masses; the density
dependence is relevant for material media and for the two-fluid
cosmological structure.\ This equation of state guarantees that $\Ka$ is
constant under compression, which is experimentally confirmed by:
%
\begin{itemize}[leftmargin=2em]
  \item Velocity time dilation measurements spanning $10^{-6}\,c$ to $0.9994\,c$
    (precision $\sim 10^{-8}$).
  \item The Fresnel drag coefficient $1 - 1/n^2$, exact across materials with
    density ratios up to $\sim$6:1.
\end{itemize}


% --------------------------------------------------
\subsection{Gravity}
% --------------------------------------------------

Gravity arises from the residual asymmetry of transient bind--unbind cycling.
%
\paragraph{Inflow velocity.}
%
\begin{equation}
  \vin(r) = \sqrt{\frac{2GM}{r}}
  \label{eq:sum-inflow}
\end{equation}
%
\paragraph{Density perturbation from gravity.}
Bernoulli's equation for free-fall inflow gives $\Delta\rho/\rhoa = 0$
exactly: the kinetic energy of the inflow balances the gravitational
potential, leaving no energy for compression.
%
\paragraph{Gravitational time dilation (from inflow velocity).}
%
\begin{equation}
  \frac{\Delta t}{t} = \frac{\vin^2}{2c^2} = \frac{GM}{r\,c^2}
  \label{eq:sum-grav-time-dilation}
\end{equation}
%
A static observer in the inflowing Aether experiences velocity time
dilation (Painlev\'{e}--Gullstrand result).

\paragraph{Event horizon (sonic horizon).}
%
\begin{equation}
  r_s = \frac{2GM}{c^2}, \qquad \vin(r_s) = c
  \label{eq:sum-schwarzschild}
\end{equation}

\paragraph{Gravitational lensing.}
The full bending angle from the acoustic metric of the inflowing Aether
(transverse advection plus time-delay asymmetry):
%
\begin{equation}
  \theta = \frac{4GM}{r_0\,c^2}
  \label{eq:sum-lensing}
\end{equation}

\paragraph{Force constant hierarchy.}
%
\begin{equation}
  G = \epsilon\,\frac{k\,(\text{geometric and mass factors})}{c^2}
  \label{eq:sum-G-from-k}
\end{equation}
%
where $\epsilon \ll 1$ is the fractional asymmetry of the bind--unbind cycle,
encoding the electromagnetic-to-gravitational force ratio between two
electrons, $ke^2/(Gm_e^2) \approx 4.17 \times 10^{42}$.

\paragraph{Hawking temperature.}
From the inflow velocity gradient at the sonic horizon:
%
\begin{equation}
  T_H = \frac{\hbar\,c^3}{8\pi\,G\,M\,k_B}
  \label{eq:sum-hawking}
\end{equation}


% --------------------------------------------------

\section*{Acknowledgments}
\addcontentsline{toc}{section}{Acknowledgments}

The author acknowledges the use of Claude (Anthropic) for \LaTeX{} formatting
and editorial assistance.

\bigskip


% ============================================================
% BIBLIOGRAPHY
% ============================================================

\begin{thebibliography}{99}

\bibitem{wilson2011}
Wilson, C.~M., Johansson, G., Pourkabirian, A., et al.,
``Observation of the dynamical Casimir effect in a superconducting circuit,''
\textit{Nature}, \textbf{479}, 376--379 (2011).



\bibitem{eddington1920}
Dyson, F.~W., Eddington, A.~S., and Davidson, C.,
``A Determination of the Deflection of Light by the Sun's Gravitational Field, from Observations Made at the Total Eclipse of May 29, 1919,''
\textit{Phil.\ Trans.\ Royal Society A}, \textbf{220}, 291--333 (1920).


\bibitem{everitt2011}
Everitt, C.~W.~F., DeBra, D.~B., Parkinson, B.~W., et al.,
``Gravity Probe B: Final Results of a Space Experiment to Test General Relativity,''
\textit{Physical Review Letters}, \textbf{106}, 221101 (2011).


\bibitem{ligo2016}
Abbott, B.~P.\ et al.\ (LIGO Scientific Collaboration and Virgo Collaboration),
``Observation of Gravitational Waves from a Binary Black Hole Merger,''
\textit{Physical Review Letters}, \textbf{116}, 061102 (2016).


\bibitem{heisenberg1927}
Heisenberg, W.,
``\"{U}ber den anschaulichen Inhalt der quantentheoretischen Kinematik und Mechanik,''
\textit{Zeitschrift f\"{u}r Physik}, \textbf{43}, 172--198 (1927).


\bibitem{svancara2024}
Svancara, P., Smaniotto, P., Solidoro, L., et al.,
``Rotating curved spacetime signatures from a giant quantum vortex,''
\textit{Nature} (2024). DOI: 10.1038/s41586-024-07176-8.


\bibitem{hathaway2021}
Hathaway, G.~D.\ and Williams, L.~L.,
``Null-result test for effect on weight from large electrostatic charge,''
\textit{Physics}, \textbf{3}, 160--172 (2021).


\bibitem{tajmar2024}
Tajmar, M., K\"ossling, M., and Neunzig, O.,
``In-depth experimental search for a coupling between gravity and electromagnetism with steady fields,''
\textit{Sci.\ Rep.}, \textbf{14}, 19427 (2024).



\bibitem{Ciufolini2004}
I.~Ciufolini and E.\,C.~Pavlis,
``A confirmation of the general relativistic prediction of the
Lense--Thirring effect,''
\emph{Nature}\ \textbf{431}, 958--960 (2004).



\bibitem{bertotti2003}
Bertotti, B., Iess, L., and Tortora, P.,
``A test of general relativity using radio links with the Cassini spacecraft,''
\textit{Nature}, \textbf{425}, 374--376 (2003).

\bibitem{gw170817}
Abbott, B.~P.\ et al.\ (LIGO Scientific Collaboration and Virgo Collaboration),
``Gravitational Waves and Gamma-Rays from a Binary Neutron Star Merger: GW170817 and GRB~170817A,''
\textit{Astrophysical Journal Letters}, \textbf{848}, L13 (2017).

\bibitem{otto-overview}
Otto, E.,
``Unifying Theory of Dynamic Aether Mechanics: General Overview,''
(2026).

\bibitem{otto-cosmology}
Otto, E.,
``Unifying Theory of Dynamic Aether Mechanics: Cosmology,''
(2026).

\bibitem{otto-em}
Otto, E.,
``Unifying Theory of Dynamic Aether Mechanics: Magnetic and Electric Fields,''
(2026).
\end{thebibliography}

\end{document}